\documentclass[11pt]{article}

\usepackage[letterpaper,margin=1in]{geometry}
\usepackage[T1]{fontenc}
\usepackage{amsmath,amsthm,mathtools}
\IfFileExists{newtxtext.sty}{\usepackage{newtxtext,newtxmath}}{\usepackage{mathptmx}\usepackage{amssymb}}
\usepackage{microtype}
\usepackage{setspace}
\usepackage{booktabs,threeparttable,array,tabularx}
\usepackage{graphicx}
\usepackage{enumitem}
\usepackage{caption}
\usepackage{placeins}
\usepackage{natbib}
\usepackage{xcolor}
\usepackage[hidelinks]{hyperref}
\usepackage[nameinlink,noabbrev]{cleveref}

\graphicspath{{../../figures/}}
\setstretch{1.5}
\setlength{\emergencystretch}{1em}
\setlength{\parskip}{0.25em}
\setlength{\parindent}{1.25em}
\captionsetup{font=small,labelfont=bf}
\setlist[itemize]{leftmargin=1.4em,itemsep=0.2em,topsep=0.3em}
\setlist[enumerate]{leftmargin=1.6em,itemsep=0.2em,topsep=0.3em}

\newtheorem{proposition}{Proposition}
\newtheorem{theorem}{Theorem}
\newtheorem{lemma}{Lemma}
\newtheorem{corollary}{Corollary}
\theoremstyle{definition}
\newtheorem{definition}{Definition}
\theoremstyle{remark}
\newtheorem{remark}{Remark}

\newcommand{\E}{\mathbb{E}}
\newcommand{\Prob}{\mathbb{P}}
\newcommand{\1}{\mathbf{1}}
\newcommand{\F}{\mathcal{F}}
\newcommand{\G}{\mathcal{G}}
\newcommand{\V}{\mathcal{V}}
\newcommand{\Loss}{\mathcal{L}}
\newcommand{\opt}{\mathrm{opt}}
\newcommand{\priv}{\mathrm{private}}
\newcommand{\cons}{\mathrm{consensus}}
\newcommand{\mkt}{\mathrm{market}}
\newcommand{\plan}{\mathrm{planner}}
\newcommand{\PoA}{\mathrm{PoA}}
\DeclareMathOperator*{\argmax}{arg\,max}

\makeatletter
% Generated from the journal main manuscript; refresh with build.py.
\newlabel{sec:benchmark}{{2}{3}{The sixteen-box benchmark}{section.2}{}}
\newlabel{sec:benchmark@cref}{{[section][2][]2}{[1][3][]3}}
\newlabel{eq:private}{{1}{4}{Private clue-following}{equation.2.1}{}}
\newlabel{eq:private@cref}{{[equation][1][]1}{[1][3][]4}}
\newlabel{eq:consensus}{{2}{4}{Shared consensus}{equation.2.2}{}}
\newlabel{eq:consensus@cref}{{[equation][2][]2}{[1][4][]4}}
\newlabel{eq:planner}{{3}{4}{Shared portfolio}{equation.2.3}{}}
\newlabel{eq:planner@cref}{{[equation][3][]3}{[1][4][]4}}
\newlabel{eq:redundancy}{{4}{5}{Shared portfolio}{equation.2.4}{}}
\newlabel{eq:redundancy@cref}{{[equation][4][]4}{[1][5][]5}}
\newlabel{fig:frontier}{{1}{6}{The exact pooled-information frontier for the canonical benchmark. The one-action consensus value is only the left endpoint. Seven coordinated actions recover the performance of eight decentralized clue-followers.\relax }{figure.caption.5}{}}
\newlabel{fig:frontier@cref}{{[figure][1][]1}{[1][6][]6}}
\newlabel{eq:distinct-private}{{5}{6}{The full action-budget frontier}{equation.2.5}{}}
\newlabel{eq:distinct-private@cref}{{[equation][5][]5}{[1][6][]6}}
\newlabel{sec:framework}{{3}{7}{Information value and protocol loss}{section.3}{}}
\newlabel{sec:framework@cref}{{[section][3][]3}{[1][7][]7}}
\newlabel{eq:frontier-general}{{6}{7}{Information value and protocol loss}{equation.3.6}{}}
\newlabel{eq:frontier-general@cref}{{[equation][6][]6}{[1][7][]7}}
\newlabel{eq:protocol-loss}{{7}{7}{Information value and protocol loss}{equation.3.7}{}}
\newlabel{eq:protocol-loss@cref}{{[equation][7][]7}{[1][7][]7}}
\newlabel{prop:voi}{{1}{7}{Value of information}{proposition.1}{}}
\newlabel{prop:voi@cref}{{[proposition][1][]1}{[1][7][]7}}
\newlabel{sec:signals}{{3.1}{8}{Independent symmetric signals}{subsection.3.1}{}}
\newlabel{sec:signals@cref}{{[subsection][1][3]3.1}{[1][8][]8}}
\newlabel{eq:posterior-count}{{8}{8}{Independent symmetric signals}{equation.3.8}{}}
\newlabel{eq:posterior-count@cref}{{[equation][8][]8}{[1][8][]8}}
\newlabel{rem:conformity}{{1}{8}{Partial conformity is monotonically harmful}{remark.1}{}}
\newlabel{rem:conformity@cref}{{[remark][1][]1}{[1][8][]8}}
\newlabel{sec:selfish}{{4}{9}{Selfish search and the price of coordination}{section.4}{}}
\newlabel{sec:selfish@cref}{{[section][4][]4}{[1][9][]9}}
\newlabel{eq:reciprocal}{{9}{10}{Payoffs, potential, and total welfare}{equation.4.9}{}}
\newlabel{eq:reciprocal@cref}{{[equation][9][]9}{[1][9][]10}}
\newlabel{eq:phi}{{10}{10}{Payoffs, potential, and total welfare}{equation.4.10}{}}
\newlabel{eq:phi@cref}{{[equation][10][]10}{[1][10][]10}}
\newlabel{eq:payoff}{{11}{10}{Payoffs, potential, and total welfare}{equation.4.11}{}}
\newlabel{eq:payoff@cref}{{[equation][11][]11}{[1][10][]10}}
\newlabel{eq:harmonic-potential}{{12}{10}{Potential and payout identity}{equation.4.12}{}}
\newlabel{eq:harmonic-potential@cref}{{[equation][12][]12}{[1][10][]10}}
\newlabel{prop:waterfill}{{5}{11}{Anonymous symmetric equilibrium}{proposition.5}{}}
\newlabel{prop:waterfill@cref}{{[proposition][5][]5}{[1][11][]11}}
\newlabel{eq:waterfill}{{13}{11}{Anonymous symmetric equilibrium}{equation.4.13}{}}
\newlabel{eq:waterfill@cref}{{[equation][13][]13}{[1][11][]11}}
\newlabel{prop:welfare}{{6}{11}{Welfare identity}{proposition.6}{}}
\newlabel{prop:welfare@cref}{{[proposition][6][]6}{[1][11][]11}}
\newlabel{eq:welfare-identity}{{14}{11}{Welfare identity}{equation.4.14}{}}
\newlabel{eq:welfare-identity@cref}{{[equation][14][]14}{[1][11][]11}}
\newlabel{thm:poa}{{1}{12}{Exact mixed price of anarchy}{theorem.1}{}}
\newlabel{thm:poa@cref}{{[theorem][1][]1}{[1][12][]12}}
\newlabel{eq:poa}{{15}{12}{Exact mixed price of anarchy}{equation.4.15}{}}
\newlabel{eq:poa@cref}{{[equation][15][]15}{[1][12][]12}}
\newlabel{eq:smooth-step}{{16}{12}{The exact price of anarchy}{equation.4.16}{}}
\newlabel{eq:smooth-step@cref}{{[equation][16][]16}{[1][12][]12}}
\newlabel{eq:uniform}{{17}{13}{Uniform posterior}{equation.4.17}{}}
\newlabel{eq:uniform@cref}{{[equation][17][]17}{[1][13][]13}}
\newlabel{eq:econstant}{{18}{13}{Uniform posterior}{equation.4.18}{}}
\newlabel{eq:econstant@cref}{{[equation][18][]18}{[1][13][]13}}
\newlabel{fig:four-protocols}{{2}{14}{Exact ex-ante discovery under four protocols in the independent sixteen-box benchmark. The labels inside the bars report the expected number of distinct boxes searched. The dashed line marks the blind coordinated portfolio at $8/16=0.50$: informed consensus falls below a protocol that ignores the clues entirely. Equal sharing supplies partial dispersion, but not enough to match natural coverage from independent private clues.\relax }{figure.caption.8}{}}
\newlabel{fig:four-protocols@cref}{{[figure][2][]2}{[1][14][]14}}
\newlabel{eq:canonical-market}{{19}{14}{The canonical four-protocol ordering}{equation.4.19}{}}
\newlabel{eq:canonical-market@cref}{{[equation][19][]19}{[1][13][]14}}
\newlabel{tab:canonical}{{1}{15}{Canonical protocol comparison: $M=16$, $N=8$, $p=0.20$.\relax }{table.caption.9}{}}
\newlabel{tab:canonical@cref}{{[table][1][]1}{[1][14][]15}}
\newlabel{sec:mechanism}{{5}{15}{Reward the rescue, not the win}{section.5}{}}
\newlabel{sec:mechanism@cref}{{[section][5][]5}{[1][15][]15}}
\newlabel{sec:copying}{{6}{16}{Common-source dependence}{section.6}{}}
\newlabel{sec:copying@cref}{{[section][6][]6}{[1][16][]16}}
\newlabel{eq:effective-channels}{{20}{17}{Channel collapse}{equation.6.20}{}}
\newlabel{eq:effective-channels@cref}{{[equation][20][]20}{[1][17][]17}}
\newlabel{eq:private-copying}{{21}{17}{Channel collapse}{equation.6.21}{}}
\newlabel{eq:private-copying@cref}{{[equation][21][]21}{[1][17][]17}}
\newlabel{eq:planner-gain}{{22}{18}{The centralized planner gain}{equation.6.22}{}}
\newlabel{eq:planner-gain@cref}{{[equation][22][]22}{[1][18][]18}}
\newlabel{eq:planner-limit}{{23}{18}{Monotone planner gain}{equation.6.23}{}}
\newlabel{eq:planner-limit@cref}{{[equation][23][]23}{[1][18][]18}}
\newlabel{eq:fill-identity}{{24}{18}{The centralized planner gain}{equation.6.24}{}}
\newlabel{eq:fill-identity@cref}{{[equation][24][]24}{[1][18][]18}}
\newlabel{fig:planner-gain}{{3}{19}{Exact canonical planner gain $V_8(c)-G_{\priv }(c)$. As common-source copying creates more duplicated capacity, identifying and reallocating it becomes more valuable.\relax }{figure.caption.11}{}}
\newlabel{fig:planner-gain@cref}{{[figure][3][]3}{[1][18][]19}}
\newlabel{tab:copying}{{2}{19}{Canonical outcomes under common-cue copying.\relax }{table.caption.12}{}}
\newlabel{tab:copying@cref}{{[table][2][]2}{[1][19][]19}}
\newlabel{fig:copying-protocols}{{4}{20}{Exact canonical comparison under common-source copying. The symmetric equal-split market initially trails decentralized report-following, then overtakes it when channel collapse becomes sufficiently severe. The crossover is a canonical computational result, not a general theorem.\relax }{figure.caption.13}{}}
\newlabel{fig:copying-protocols@cref}{{[figure][4][]4}{[1][20][]20}}
\newlabel{eq:crossover}{{25}{20}{Selfish search under copying}{equation.6.25}{}}
\newlabel{eq:crossover@cref}{{[equation][25][]25}{[1][19][]20}}
\newlabel{sec:implications}{{8}{23}{Organizational design implications}{section.8}{}}
\newlabel{sec:implications@cref}{{[section][8][]8}{[1][23][]23}}
\newlabel{sec:limitations}{{9}{24}{Limitations and extensions}{section.9}{}}
\newlabel{sec:limitations@cref}{{[section][9][]9}{[1][24][]24}}

\makeatother
\renewcommand{\thesection}{EC.\arabic{section}}
\renewcommand{\theequation}{EC.\arabic{equation}}
\renewcommand{\theproposition}{EC.\arabic{proposition}}
\renewcommand{\thelemma}{EC.\arabic{lemma}}
\renewcommand{\theremark}{EC.\arabic{remark}}
\hypersetup{draft}
\hypersetup{pdfauthor={},pdftitle={Electronic Companion: The Shared Discovery Paradox}}
\title{Electronic Companion: The Shared Discovery Paradox}
\author{}
\date{}
\begin{document}
\maketitle
This companion uses the main manuscript's notation and its nondegenerate domain $M\ge2$, $2\le N\le M$, $1/M<p<1$, unless a boundary case is stated explicitly. Section numbers EC.1--EC.6 refer to this document; other labels refer to the main manuscript.



\section{Exact independent enumeration}
\label{app:enumeration}

By symmetry, condition on box $1$ being the target. The report-count vector
\[
C=(C_1,\ldots,C_M)\sim\mathrm{Multinomial}(N;p,q,\ldots,q).
\]
For a count vector $x$, let
\[
r(x)=|\{b:x_b>x_1\}|,
\qquad
t(x)=|\{b:x_b=x_1\}|.
\]
With uniform randomization at cutoff ties, the conditional probability that the true box belongs to a top-$L$ count set is
\begin{equation}
\tau_L(x)=
\begin{cases}
0, & r(x)\ge L,\\
1, & r(x)+t(x)\le L,\\
\dfrac{L-r(x)}{t(x)}, & r(x)<L<r(x)+t(x).
\end{cases}
\label{eq:tie-inclusion}
\end{equation}
Hence the exact pooled frontier is
\begin{equation}
G_{M,N,L}(p)
=\sum_{x_1+\cdots+x_M=N}
\frac{N!}{\prod_{b=1}^M x_b!}
 p^{x_1}q^{N-x_1}\tau_L(x).
\label{eq:multinomial-frontier}
\end{equation}
For $M=16$ and $N=8$, the sum contains
\[
\binom{N+M-1}{M-1}=\binom{23}{15}=490{,}314
\]
count vectors. Grouping by permutations of the fifteen non-target boxes reduces both the independent and common-cue calculations to $67$ count classes.

At $L=7$, the difference between the pooled portfolio and private clue-following can be decomposed as
\[
\Prob(\theta\in S^*_7\setminus S_{\priv})
-\Prob(\theta\in S_{\priv}\setminus S^*_7).
\]
Exact enumeration gives a synthesis gain of $0.011744$ and a compression loss of $0.007962$. At $L=8$, no box named by a private clue is discarded; the improvement $0.027193$ comes entirely from reallocating duplicated capacity.

\section{Common-cue likelihood and Bayes ranking}
\label{app:likelihood}

For a candidate target $t$, define the report probability vector
\[
\rho_r^{(t)}=
\begin{cases}
p,&r=t,\\q,&r\neq t.
\end{cases}
\]
Let $z$ be the common cue and $k$ the copier count. For observed counts $x$,
\begin{align}
&\Prob(C=x\mid\theta=t,X_0=z,K=k)\\
&\quad=
\1\{x_z\ge k\}
\frac{(N-k)!}{(x_z-k)!\prod_{r\neq z}x_r!}
\prod_{r=1}^M\left(\rho_r^{(t)}\right)^{x_r-k\1\{r=z\}}.
\end{align}
Marginalizing gives
\begin{align}
\Prob(C=x\mid\theta=t,c)
=\sum_{z=1}^M\rho_z^{(t)}
\sum_{k=0}^N\binom Nk c^k(1-c)^{N-k}
\Prob(C=x\mid\theta=t,X_0=z,K=k).
\label{eq:common-likelihood}
\end{align}
With a uniform prior, normalize these $M$ likelihoods to obtain the herd-aware posterior. The verification script precomputes the coefficient of every basis term $c^k(1-c)^{N-k}$ for each count class and candidate target.

\section{Proof of the monotone planner-gain theorem}
\label{app:planner-proof}

We give the omitted details for \cref{eq:fill-identity} and monotonicity.

\begin{lemma}[Named states precede unnamed states]
For every realizable report vector under the common-cue model, every state with a positive report count has posterior probability at least as large as every state with zero reports.
\end{lemma}

\begin{proof}
Fix a named state $b$ and an unnamed state $d$. Compare the mixture likelihood under hypotheses $\theta=b$ and $\theta=d$ term by term in \cref{eq:common-likelihood}. If the common cue is neither $b$ nor $d$, the cue prefactor is identical while the residual multinomial likelihood differs by a nonnegative power of $p/q$ in favor of $b$. If the cue is $b$, its prefactor is $p$ under $\theta=b$ and $q$ under $\theta=d$, and the residual likelihood contributes another nonnegative power of $p/q$. If the cue is $d$, feasibility requires $k=0$ because $d$ is unnamed; the combined likelihood ratio is $(p/q)^{x_b-1}\ge1$. Every matched nonzero term weakly favors the named state. Summing over cues and copier counts preserves the inequality.
\end{proof}

The optimal $N$-action portfolio therefore includes all $D$ distinct named states and fills $N-D$ positions with unnamed states. On the event that private report-following misses, the target lies among the $M-D$ unnamed states. Symmetry assigns equal conditional posterior mass to those states, so the planner's conditional rescue probability is $(N-D)/(M-D)$. When $D=M$, the miss event is impossible; the product in \cref{eq:fill-identity} is defined as zero rather than evaluating $0/0$. This proves \cref{eq:fill-identity}.

Condition on copier count $K=k$. Given $k$, the joint law of the common cue and the $N-k$ independent clues does not otherwise depend on $c$. Define
\[
h(k)=\E\left[\left.\1\{\text{miss}\}\frac{N-D}{M-D}\right|K=k\right].
\]
Then $\Delta_{\plan}(c)=\E[h(K)]$ for $K\sim\mathrm{Binomial}(N,c)$. We have $h(0)=h(1)$ because one copy of $X_0$ has the same joint distribution as one independent clue. For $k\ge1$, couple the $k$-copier system and the $(k+1)$-copier system by converting one independent searcher into an additional copier while holding the common cue and all other reports fixed. Because the common cue is already represented, the conversion can only weakly reduce $D$. It can only weakly increase the miss indicator. The fraction $(N-D)/(M-D)$ is nondecreasing as $D$ falls when $N\le M$. Thus $h(k+1)\ge h(k)$. For $1/M<p<1$ and $N\ge2$, the inequality is strict at $k=N-1$: with positive probability, the converted independent clue names the target (probability $p>0$) while the already represented common cue misses (probability $1-p>0$). Before conversion, the miss integrand is zero; after conversion, $D=1$ and it equals $(N-1)/(M-1)>0$. This event has positive probability $p(1-p)$, giving $h(N)>h(N-1)$.

If $K'\sim\mathrm{Binomial}(N-1,c)$, the standard binomial derivative identity yields
\[
\frac{d}{dc}\E[h(K)]
=N\E[h(K'+1)-h(K')]>0
\qquad(c\in(0,1)).
\]
At $p=1$, both protocols succeed surely and the gain is zero, so strictness requires the stated upper bound on $p$. The hypothesis $N\ge2$ is also necessary: with a single searcher the named box is the posterior mode, the planner replicates private search, the gain is identically zero, and the derivative vanishes throughout $(0,1)$ because $h(0)=h(1)$.
At $c=1$, one box is named. The planner adds $N-1$ of the $M-1$ unnamed boxes, each with posterior mass $q$, while private search covers only the named box. Hence $\Delta_{\plan}(1)=(N-1)q$.

\section{Proof of the proportional sparse-evidence limit}
\label{app:scaling}

Throughout, $N=N_M$ with $N_M/M\to\alpha\in(0,1)$, and $p_M=r/(M-1+r)$, $q_M=1/(M-1+r)$ for fixed $r>1$, so that $N_Mp_M\to\alpha r$ and $N_Mq_M\to\alpha$: the target receives an asymptotically $\mathrm{Poisson}(\alpha r)$ number of reports, while each false box receives approximately $\mathrm{Poisson}(\alpha)$. All limits are as $M\to\infty$ and all clues are independent.

\emph{Blind portfolio.} The coordinator selects $N_M$ distinct boxes without observing clues. Under the uniform prior, $G_{\mathrm{blind}}=N_M/M\to\alpha$.

\emph{Private clue-following.} By \cref{eq:private}, $G_{\priv}=1-(1-p_M)^{N_M}$, and $N_M\log(1-p_M)=-N_Mp_M\bigl(1+O(p_M)\bigr)\to-\alpha r$, so $G_{\priv}\to1-e^{-\alpha r}$.

\emph{Consensus.} Condition on $\theta$. The target count $C_\theta\sim\mathrm{Binomial}(N_M,p_M)$ converges in distribution to $\mathrm{Poisson}(\alpha r)$ and is therefore bounded in probability. Conditional on the number of off-target reports $F=N_M-C_\theta$, the false counts are multinomial with $F$ trials over $M-1$ equally likely boxes, and $F/M\to\alpha$ in probability. Fix an integer $j$ and let $W_j=|\{b\neq\theta:C_b\ge j\}|$. Standard occupancy asymptotics give $\E[W_j]/M\to\Prob(\mathrm{Poisson}(\alpha)\ge j)>0$, and the variance of $W_j$ is $O(M)$ because the occupancy indicators are negatively correlated, so $W_j\to\infty$ in probability. Hence $\max_{b\neq\theta}C_b\ge j$ with probability tending to one for every fixed $j$: the largest false count diverges in probability while $C_\theta$ remains tight. Consensus success requires $C_\theta\ge\max_{b\neq\theta}C_b$ (an event whose probability tends to zero), so $G_{\cons}\to0$; uniform tie-breaking only lowers the success probability further.

\emph{Portfolio.} By \cref{eq:posterior-count}, every named box has strictly larger posterior mass than every unnamed box, and at most $N_M$ boxes are named, so the top-$N_M$ rule retains every named box and fills its remaining slots uniformly at random among the posterior-tied unnamed boxes. Let $D$ be the number of distinct named boxes. On $\{C_\theta\ge1\}$ the portfolio contains $\theta$. On the miss event $\{C_\theta=0\}$, the target is one of the $M-D$ unnamed boxes, which carry equal conditional posterior mass, so the conditional rescue probability is $(N_M-D)/(M-D)$. This is the $c=0$ case of \cref{eq:fill-identity}:
\[
G_{\mathrm{portfolio}}
=G_{\priv}+\E\!\left[\1\{C_\theta=0\}\,\frac{N_M-D}{M-D}\right].
\]
Conditional on the miss event, $D$ is the number of occupied boxes when $N_M$ balls are placed uniformly over $M-1$ boxes. Occupancy concentration ($\E[D]=(M-1)[1-(1-\tfrac{1}{M-1})^{N_M}]$ with variance $O(M)$) gives $D/M\to1-e^{-\alpha}$ in probability, hence
\[
\frac{N_M-D}{M-D}
\;\longrightarrow\;
\tau_\alpha
=\frac{\alpha-(1-e^{-\alpha})}{e^{-\alpha}}
=1-(1-\alpha)e^{\alpha}
\qquad\text{in probability},
\]
where $\tau_\alpha\in(0,1)$ because $e^{-\alpha}>1-\alpha$ for $\alpha>0$: duplicated reports leave free slots asymptotically. The ratio is bounded by one, so bounded convergence together with $\Prob(C_\theta=0)\to e^{-\alpha r}$ yields
\[
G_{\mathrm{portfolio}}
\longrightarrow
\bigl(1-e^{-\alpha r}\bigr)+e^{-\alpha r}\bigl[1-(1-\alpha)e^{\alpha}\bigr]
=1-(1-\alpha)\,e^{-\alpha(r-1)}.
\]
Finally, $(1-\alpha)e^{\alpha}$ is strictly decreasing on $(0,1)$ from the value one at $\alpha=0$, so $(1-\alpha)e^{\alpha}<1$ and the portfolio limit strictly exceeds the private limit for every $\alpha\in(0,1)$. \qed

\emph{Symmetric equal-split market.} Write the scaled posterior weight of box $b$ as $W_b=M\pi_b=r^{C_b}/(Z_M/M)$ with $Z_M=\sum_b r^{C_b}$, and recall from \cref{prop:scaling} the notation $f_j=e^{-\alpha}\alpha^j/j!$, $\psi_j=r^je^{-\alpha(r-1)}$, $g_j=f_j\psi_j=e^{-\alpha r}(\alpha r)^j/j!$, and $h(x)=(1-e^{-x})/x$.

\emph{Step 1: weighted occupancy law of large numbers.} For a false box, $\E\bigl[r^{C_b}\bigr]=\bigl(1+(r-1)q_M\bigr)^{N_M}\to e^{\alpha(r-1)}$ and $\E\bigl[r^{2C_b}\bigr]=\bigl(1+(r^2-1)q_M\bigr)^{N_M}\to e^{\alpha(r^2-1)}<\infty$. For two false boxes $a\neq b$, each clue lands on at most one of them, so $\E\bigl[r^{C_a}r^{C_b}\bigr]=\bigl(1+2(r-1)q_M\bigr)^{N_M}$, and a direct expansion shows the covariance is $O(1/M)$ times a bounded factor. Chebyshev's inequality then gives $Z_M/M\to e^{\alpha(r-1)}$ in probability, so $W_b\to\psi_j$ in probability for a box with $C_b=j$; the target term has bounded expectation $\E[r^{C_\theta}]=(1+(r-1)p_M)^{N_M}=O(1)$ and hence contributes $o_{\Prob}(1)$ after division by $M$. For false-box tails, $r^{C_b}\1\{C_b\ge J\}\le r^{-J}r^{2C_b}$, so their normalized tail expectation is uniformly $O(r^{-J})$; Markov's inequality makes the tails uniformly small in probability. Finite count-class frequencies converge by occupancy concentration; truncation then yields, for the weighted sums, $M^{-1}\sum_b r^{C_b}\1\{C_b\ge J\}\to e^{\alpha(r-1)}\sum_{j\ge J}g_j$ in probability, with the tail uniformly small in $J$ because $\sum_j g_j=1$ (the identity $f_j\psi_j=g_j$ is the exponential tilting of a Poisson law).

\emph{Step 2: convergence of the budget equation.} The same tail bound gives $\max_b W_b/M\to0$ in probability; thus for any fixed $\Lambda>0$ the lower inverse endpoint $\max_b W_b/N_M$ tends to zero. The finite inverse/budget equation is therefore in its non-pure domain with probability tending to one. At the finite anonymous equilibrium of \cref{prop:waterfill}, an active box satisfies $W_b\,\phi(s_b)=\Lambda_M:=M\lambda_M$, and the budget constraint $\sum_b s_b=1$ reads $M^{-1}\sum_b x_b=N_M/M$ with $x_b=N_Ms_b$. Since $(1-x/N)^N\to e^{-x}$ uniformly on compact $x$-sets, $\phi(x/N)\to h(x)$ uniformly on compacts, with $h$ continuous and strictly decreasing from $h(0)=1$. Define the limit budget function
\[
B(\Lambda)=\sum_{j\,:\,\psi_j>\Lambda} f_j\,h^{-1}\!\bigl(\Lambda/\psi_j\bigr).
\]
Because $h(x)\le 1/x$, every active class obeys $x_j\le\psi_j/\Lambda$, so $B(\Lambda)\le\Lambda^{-1}\sum_jg_j=1/\Lambda$ and the series converges locally uniformly; $B$ is therefore continuous, strictly decreasing where positive, with $B(0^+)=\infty$ and $B(\infty)=0$, and \cref{eq:market-limit} has a unique root $\Lambda$. The same bound $x_b\le W_b/\Lambda$ controls the high-count contribution to the finite budget uniformly in $M$ by Step 1, so the finite (random) budget functions converge to $B$ in probability at every fixed $\Lambda$. Strict monotonicity then squeezes the equilibrium value: for any $\Lambda'<\Lambda<\Lambda''$, eventually $B_M(\Lambda')>N_M/M>B_M(\Lambda'')$ with probability tending to one, so $\Lambda_M\to\Lambda$ in probability.

\emph{Step 3: from values to discovery.} The welfare identity of \cref{prop:welfare} gives conditional discovery $N_M\lambda_M=(N_M/M)\Lambda_M\le1$, so $\Lambda_M\le M/N_M$ is uniformly bounded, and bounded convergence yields $G_{\mkt}^{\mathrm{sym}}=\E[N_M\lambda_M]\to\alpha\Lambda$.

\emph{Step 4: the tilted form.} Multiplying the budget equation by $\Lambda$ and using $\Lambda=\psi_jh(x_j)$ on active classes,
\[
\alpha\Lambda
=\sum_j f_jx_j\,\psi_jh(x_j)
=\sum_j f_j\psi_j\bigl(1-e^{-x_j}\bigr)
=\sum_j g_j\bigl(1-e^{-x_j}\bigr):
\]
in the limit the market covers the target's own count class $j$ with probability $1-e^{-x_j}$, where the target's count is $\mathrm{Poisson}(\alpha r)$ and the number of occupants of any fixed box is asymptotically $\mathrm{Poisson}(x_j)$. Canonically $\Lambda=1.094021$ lies strictly between $\psi_1=0.948$ and $\psi_2=3.555$, so the limit market abandons every box with fewer than two reports. \qed

\emph{Abundant evidence (\cref{rem:abundant}).} Hold $p$ fixed. By the multiplicative Chernoff bound, $\Prob(C_\theta\le N_Mp/2)\le e^{-N_Mp/8}\to0$. Each false count is $\mathrm{Binomial}(N_M,q_M)$ with mean $N_Mq_M\to\alpha(1-p)$, so the standard upper-tail bound gives $\Prob(C_b\ge N_Mp/2)\le\bigl(2eq_M/p\bigr)^{N_Mp/2}$, whose base is $O(1/M)$ while the exponent grows linearly in $M$; the union bound over the $M-1$ false boxes therefore vanishes. With probability tending to one the target is the unique plurality winner, so $G_{\cons}\to1$. The same argument applies whenever $\lambda_M=N_Mp_M$ satisfies $\lambda_M/\log M\to\infty$: the target count concentrates above $\lambda_M/2$ while $M\exp\bigl[-(\lambda_M/2)\log(\lambda_M/2e\alpha)\bigr]\to0$. \qed

\section{The proportional sparse-evidence limit}
\label{sec:scaling}

The sixteen-box environment is deliberately small, so it is natural to ask whether the ordering is a small-number artifact. Scale the environment proportionally: let the number of searchers grow with the number of boxes, $N_M/M\to\alpha\in(0,1)$, and hold the clue likelihood ratio $r=p_M/q_M>1$ fixed, so that
\[
p_M=\frac{r}{M-1+r},
\qquad
q_M=\frac{1}{M-1+r}.
\]
The canonical benchmark corresponds to $\alpha=1/2$ and $r=3.75$. Fixing $r$ is one natural invariant among several---one could instead preserve mutual information---but it keeps the report-count occupancy problem nondegenerate: the target receives an asymptotically $\mathrm{Poisson}(\alpha r)$ number of reports, while a typical false box receives approximately $\mathrm{Poisson}(\alpha)$.

\begin{proposition}[Proportional sparse-evidence limit]
\label{prop:scaling}
In the independent symmetric model with $N_M/M\to\alpha\in(0,1)$ and fixed likelihood ratio $r>1$, consider the blind coordinated portfolio ($N_M$ distinct boxes chosen without clues), pooled one-action consensus, private clue-following, and the pooled top-$N_M$ portfolio. As $M\to\infty$,
\[
G_{\mathrm{blind}}\longrightarrow\alpha,
\qquad
G_{\cons}\longrightarrow 0,
\qquad
G_{\priv}\longrightarrow 1-e^{-\alpha r},
\qquad
G_{\mathrm{portfolio}}\longrightarrow 1-(1-\alpha)\,e^{-\alpha(r-1)}.
\]
The anonymous symmetric equal-split market on the pooled posterior also converges:
\begin{equation}
G_{\mkt}^{\mathrm{sym}}\longrightarrow \alpha\Lambda,
\qquad
\text{where $\Lambda$ uniquely solves}
\quad
\sum_{j\,:\,\psi_j>\Lambda} f_j\,h^{-1}\!\bigl(\Lambda/\psi_j\bigr)=\alpha,
\label{eq:market-limit}
\end{equation}
with $f_j=e^{-\alpha}\alpha^{j}/j!$, $\psi_j=r^{j}e^{-\alpha(r-1)}$, and $h(x)=(1-e^{-x})/x$ strictly decreasing from $h(0)=1$. Writing $x_j=h^{-1}(\Lambda/\psi_j)$ on active count classes and $g_j=f_j\psi_j=e^{-\alpha r}(\alpha r)^{j}/j!$, the equivalent tilted form is $G_{\mkt}^{\mathrm{sym}}\to\sum_j g_j\bigl(1-e^{-x_j}\bigr)$.
\end{proposition}

The proof is in \cref{app:scaling}. For the canonical scaling $\alpha=1/2$, $r=3.75$, the solution of \cref{eq:market-limit} is $\Lambda=1.094021$, and the complete limit ordering is
\[
\begin{gathered}G_{\cons}\to 0
\;<\;G_{\mathrm{blind}}\to 0.50
\;<\;G_{\mkt}^{\mathrm{sym}}\to 0.547011,\\
G_{\mkt}^{\mathrm{sym}}\to 0.547011
\;<\;G_{\priv}\to 0.846645
\;<\;G_{\mathrm{portfolio}}\to 0.873580.\end{gathered}
\]
Because $(1-\alpha)e^{\alpha}<1$ on $(0,1)$, the portfolio limit strictly exceeds the private limit at every $\alpha$. The market's position strictly above consensus and weakly below the portfolio also holds at every $(\alpha,r)$, since the finite inequalities pass to the limit. The positions of the blind benchmark relative to private search and to the market, by contrast, depend on $(\alpha,r)$, and the displayed ordering is a property of the canonical scaling.

The market limit is water-filling against a Poisson count profile: in the limit the equilibrium abandons every box with fewer than two reports (canonically $\psi_1=0.948<\Lambda<\psi_2=3.555$) and spreads over the high-count classes in which the target, whose own count is $\mathrm{Poisson}(\alpha r)$, is hidden. The finite canonical value $0.599099$ sits above its limit $0.547011$, and the seeded Monte Carlo diagnostic in the verification script shows a monotone decline through $M\in\{16,64,256,1024\}$ ($0.5991$, $0.5740$, $0.5582$, $0.5494$): as the posterior spreads over more comparably ranked boxes, symmetric mixing wastes more of its budget on collisions inside the leading count classes.

Two readings of the canonical limit deserve note. First, information without assignment and assignment without information deliver nearly the same discovery: the market reaches $0.547$ holding the full pooled posterior but no coordination device, the blind portfolio reaches $0.500$ holding a coordination device but no information, and only their combination reaches $0.874$. Second, the market--private gap at scale admits an exact accounting. Since $G_{\priv}\to\Prob\bigl(\mathrm{Poisson}(\alpha r)\ge1\bigr)$ and the canonical limit market abandons the zero- and one-report classes,
\[
\begin{gathered}G_{\priv}-G_{\mkt}^{\mathrm{sym}}\;\longrightarrow\;g_1+\sum_{j\ge2}g_j\,e^{-x_j}\\=0.287541+0.012094.\end{gathered}
\]
essentially the entire gap is the probability that the target receives exactly one report. The singleton discoveries that private clue-followers keep for free are precisely what an anonymous market rationally abandons; collisions inside the classes the market does cover cost only $0.012$.

\begin{remark}[Abundant-evidence regime]
\label{rem:abundant}
The paradox is a sparse-evidence phenomenon. Hold $p\in(0,1)$ fixed with $N_M/M\to\alpha\in(0,1)$, so that the likelihood ratio $r_M=p/q_M\to\infty$. Then $G_{\cons}\to1$, and a fortiori every clue-using protocol converges to one: the paradox vanishes in levels, and distinctions among diversified protocols persist only in convergence rates. The clean dial is $\lambda_M=N_Mp_M$, the target's expected report count: \cref{prop:scaling} is the bounded-$\lambda$ regime, while whenever $\lambda_M/\log M\to\infty$ the union bound in \cref{app:scaling} gives consistency. The intermediate window $\lambda_M\asymp\log M$ is a multiple-comparisons threshold problem we leave open. The benchmark is aimed accordingly: consensus institutions are adequate where evidence about the leading candidate is abundant; the compression loss binds in sparse-evidence search---frontier science, early-stage investing, novel-hypothesis generation---where the target's support never separates from the noise.
\end{remark}

The limit sharpens the finite benchmark in two ways. First, the paradox is not a small-number artifact: consensus discovery vanishes while every diversified protocol retains a strictly positive limit. Second, large-scale consensus fails through an additional channel. The target's report count remains stochastically bounded, but the maximum count among linearly many false boxes diverges, so the plurality winner is eventually an extreme false positive: top-one selection loses a multiple-comparisons contest that the portfolio, which retains every named box, never enters. The pooled ranking still contains substantial information; converting it into one repeated action is what discards it. The portfolio formula is the asymptotic form of the finite reallocation identity that reappears under copying in \cref{sec:copying}: the planner's entire limiting gain over private search, $e^{-\alpha r}\bigl[1-(1-\alpha)e^{\alpha}\bigr]$, comes from redirecting duplicated capacity to unnamed states on the event that every clue misses.

\section{Computational verification and status of claims}
\label{app:computation}

The accompanying script \texttt{verify\_shared\_discovery\_v8\_6.py} performs the following checks.

\begin{enumerate}
\item It verifies the reciprocal-share identity and the water-filling equilibrium on $200$ random posteriors, including dominant-posterior pure corners.
\item It verifies the welfare identity $G_{\mkt}^{\mathrm{sym}}=N\lambda$ posterior by posterior.
\item It reproduces the uniform-posterior formula and the $e/(e-1)$ limit.
\item It verifies the tight $2-1/N$ family and the sole-rescue pure-equilibrium characterization.
\item It exactly enumerates the $67$ canonical report-count classes for every copying value, verifies that class probabilities sum to one, and reproduces the private-search closed form.
\item It reproduces the independent pooled frontier, the canonical market value $0.599099252439$, the full copying table, and the crossover $c^*=0.788461656521$.
\item It verifies the blind-control values $8/16=0.50$ and $1-(1-1/16)^8=0.403281$, and that the blind-sentinel hybrid family is maximized at zero sentinels.
\item It computes the average-quality column exactly ($Q_{\mkt}^{\mathrm{sym}}=0.344136160273$) and verifies the redundancy identity $NQ-G=\E[(K-1)^{+}]$ class by class for the market and in closed form for the other protocols.
\item It evaluates the proportional sparse-evidence limits at $\alpha=1/2$, $r=3.75$, checks the fill-identity form of the portfolio limit, and runs a seeded Monte Carlo convergence diagnostic at $M\in\{16,64,256\}$.
\item It cross-checks a closed-form expression for the common-cue likelihood coefficients against the direct enumeration on three $(M,N,p)$ configurations, and validates a batched equilibrium solver against the scalar solver on random posteriors.
\item It solves the market limit equation \cref{eq:market-limit} to high precision ($\Lambda=1.094021073740$, $G=0.547010536870$), verifies the tilted-form identity to $10^{-11}$, confirms that the active support starts at two reports, and runs a seeded Monte Carlo convergence diagnostic for the finite market at $M\in\{16,64,256,1024\}$.
\item It computes the partial-conformity curve of \cref{rem:conformity} exactly for $a=0,\ldots,8$, verifies the endpoints $G(0)=G_{\priv}$ and $G(8)=G_{\cons}$, and confirms the strict monotone decline.
\item It verifies the limit gap decomposition $G_{\priv}-G_{\mkt}^{\mathrm{sym}}\to g_1+\sum_{j\ge2}g_je^{-x_j}$, whose two terms are $0.287541$ and $0.012094$, to $10^{-11}$, and runs a seeded fixed-$p$ Monte Carlo confirming the abundant-evidence remark (consensus $0.385\to0.947\to1.000$ at $M\in\{16,64,256\}$).
\item It scans the $72$-combination comparative-statics grid, verifies that the market--private gap changes sign at most once in every combination, and writes the crossover locations to \texttt{grid\_crossings.csv}.
\item It generates every figure and the underlying CSV tables.
\end{enumerate}

\paragraph{Analytically established.}
The value-of-information and fixed-budget dominance results; posterior-count ranking under independent symmetric signals; private coverage dominance over one-action consensus; the weak monotone harm of partial conformity by set inclusion; the blind-control values; the redundancy identity $NQ-G=\E[(K-1)^{+}]$; the proportional sparse-evidence limits, including the market cell of \cref{eq:market-limit} and the abundant-evidence consistency remark; the reciprocal-share identity; the exact-potential and payout identities for equal sharing; the anonymous water-filling equilibrium including its pure corner; $G_{\mkt}^{\mathrm{sym}}=N\lambda$; the exact mixed $\PoA=2-1/N$ upper bound and tight family; full pure-strategy implementation under sole-rescue; channel collapse and the closed-form private discovery function; failure of count ranking under dependence by explicit counterexample; the fill identity; and strict monotonicity of the centralized planner gain.

\paragraph{Exactly enumerated or numerically solved for the canonical model.}
Consensus, private, planner, and action-budget frontier values; the partial-conformity curve $G(a)$; recovery budget; expected distinct-action counts; the anonymous market value and its average-quality column; the common-cue table; and the market/private crossover. Enumeration over report classes is exact up to floating-point evaluation of finite formulas; one-dimensional equilibrium and crossover equations are solved to high numerical precision.

\paragraph{Computational patterns, not general theorems.}
The canonical herding discount vanishes for action budgets $L\ge3$ on the tested copying grid; the canonical recovery budget is nonincreasing over that grid; and the market/private gap changes sign at most once on the tested $72$-combination grid---exactly once in $52$ combinations and never in the remaining $20$, where private report-following stays ahead through $c=0.98$. These patterns should not be generalized without additional analysis.

\bibliographystyle{plainnat}
\bibliography{bibliography}
\end{document}
